\subsection{Luồng xử lý từ Public State đến verdict}
\label{sec:agent_pipeline}
\begin{sloppypar}

Tác nhân suy luận của dự án được tách thành ba tầng độc lập.
\texttt{DeductiveAgent} quyết định \emph{hỏi gì} và \emph{nộp gì};
\texttt{DPLLChecker} dựng cơ sở tri thức và trả lời câu hỏi thỏa được hay
không; \texttt{CNFEncoder} cùng hàm \texttt{solve} của bộ giải DPLL đảm nhận
phần tính toán logic. Ba tầng chỉ giao tiếp qua giao diện
\texttt{ConstraintChecker}, gồm đúng một phương thức
\begin{center}
  \texttt{has\_solution(public\_state, assumption=None) -> bool},
\end{center}
trả về \texttt{True} khi $\mathit{KB}_t$ (cùng giả thiết tạm thời nếu có) thỏa
được. Tác nhân nhận bộ kiểm tra qua hàm khởi tạo, nên cùng một mã tác nhân có
thể chạy trên \texttt{DPLLChecker} khi chơi thật hoặc trên
\texttt{BruteForceChecker} khi đối chứng trong kiểm thử. Các mục sau dùng biến
$x_i$ như ở Mục~\ref{sec:prop_vars}; literal $x_i$ ứng với ký hiệu $C_i$ (ô $i$
là tội phạm) của đề bài và của Mục~4.2.

\paragraph{Năm bước của một lượt suy luận.}
Mỗi lượt bắt đầu từ một ảnh chụp \texttt{PublicState} và kết thúc bằng một
phán quyết được Game Engine chấp nhận, hoặc bằng kết luận rằng không còn phán
quyết nào chứng minh được. Bảng~\ref{tab:agent_pipeline} nêu đầu vào và đầu
ra của từng bước.

\begin{enumerate}
  \item \textbf{Lấy thông tin công khai.} Tác nhân gọi
    \texttt{GameEngine.get\_public\_state()} và nhận ảnh chụp chỉ đọc gồm bàn
    chơi, tập lời khai đã lật $\mathcal{R}_t$ và tập phán quyết đã chứng minh
    $\mathcal{P}_t$. Ảnh chụp không chứa $\sigma^\star$ hay lời khai chưa lật
    (Mục~\ref{sec:knowledge_base}), nên mọi bước phía sau chỉ có thể dùng
    thông tin công khai.
  \item \textbf{Xác định ô cần xét.} Tác nhân duyệt bàn chơi theo thứ tự hàng
    và lấy tập ô chưa chứng minh
    $U_t = \mathcal{C} \setminus \{\, i : \langle i, s\rangle \in \mathcal{P}_t \,\}$,
    giữ thứ tự định danh tăng dần. Thứ tự cố định làm cho toàn bộ quá trình
    suy luận tất định: chạy lại cùng một bản đồ luôn cho cùng một trace.
  \item \textbf{Dựng $\mathit{KB}_t$ dưới dạng CNF.} Mỗi lần
    \texttt{has\_solution} được gọi, \texttt{DPLLChecker} tạo một
    \texttt{CNFEncoder} mới trên bàn chơi, mã hóa lần lượt ràng buộc của mọi lời
    khai trong $\mathcal{R}_t$, thêm một mệnh đề đơn vị cho mỗi phán quyết trong
    $\mathcal{P}_t$ theo \eqref{eq:kb}, rồi thêm mệnh đề đơn vị của giả thiết
    $x_i = s$ nếu có. Số biến truyền cho bộ giải là
    \texttt{next\_variable}$-1$, gồm $|\mathcal{C}|$ biến ô và các biến phụ
    trợ vừa được cấp phát. Vì công thức được dựng lại từ đầu, giả thiết của
    truy vấn trước không bao giờ rò sang truy vấn sau.
  \item \textbf{Truy vấn SAT.} Hàm \texttt{solve(cnf, n)} trả về một phép gán
    thỏa mãn hoặc \texttt{None} (Mục~4.3); \texttt{DPLLChecker} rút gọn kết quả
    thành một giá trị Boolean. Với mỗi ô $i \in U_t$, tác nhân hỏi hai câu
    $\mathrm{SAT}(\mathit{KB}_t \land x_i)$ và
    $\mathrm{SAT}(\mathit{KB}_t \land \neg x_i)$, và ghi mỗi câu hỏi vào một
    bản ghi \texttt{Query(assumption, satisfiable)}.
  \item \textbf{Phân loại và chọn verdict.} Cặp câu trả lời được phân loại
    thành \texttt{CRIMINAL}, \texttt{INNOCENT} hoặc \texttt{UNKNOWN} theo
    Mục~\ref{sec:entailment_classification}. Ô đầu tiên (theo thứ tự định danh)
    có nhãn khác \texttt{UNKNOWN} cho ra phán quyết $\langle i, s\rangle$, được
    gửi tới \texttt{GameEngine.submit\_verdict}. Engine tự kiểm tra lại bằng hai
    truy vấn của chính nó, lật lời khai của ô $i$ và trả về ảnh chụp mới
    (Mục~\ref{sec:deductive_agent}).
\end{enumerate}

\begin{table}[htbp]
  \centering\small
  \caption{Đầu vào và đầu ra của từng bước trong một lượt suy luận}
  \label{tab:agent_pipeline}
  \begin{tabularx}{\textwidth}{c>{\raggedright\arraybackslash}p{3.5cm}%
      >{\raggedright\arraybackslash}X>{\raggedright\arraybackslash}X}
    \toprule
    \textbf{Bước} & \textbf{Thành phần} & \textbf{Đầu vào} & \textbf{Đầu ra} \\
    \midrule
    1 & \texttt{GameEngine} & Trạng thái nội bộ \texttt{GameState}
      & \texttt{PublicState} chỉ đọc: bàn chơi, $\mathcal{R}_t$, $\mathcal{P}_t$ \\
    2 & \texttt{DeductiveAgent} & \texttt{PublicState}
      & Danh sách ô chưa chứng minh $U_t$, theo thứ tự định danh \\
    3 & \texttt{DPLLChecker}, \texttt{CNFEncoder}
      & \texttt{PublicState}, giả thiết \texttt{Verdict} hoặc \texttt{None}
      & CNF dạng DIMACS của $\mathit{KB}_t \land (x_i = s)$ và số biến $n$ \\
    4 & \texttt{solve} (DPLL) & CNF, $n$
      & Phép gán thỏa mãn hoặc \texttt{None}, rút gọn thành \texttt{bool} \\
    5 & \texttt{DeductiveAgent} & Hai giá trị \texttt{bool} cho mỗi ô
      & Nhãn của ô; \texttt{Move} gồm trạng thái, phán quyết và danh sách
        \texttt{Query} \\
    \bottomrule
  \end{tabularx}
\end{table}

\paragraph{Ba điểm vào của tác nhân.}
Cùng một đường ống trên được dùng qua ba phương thức.
\texttt{next\_move(state)} chạy năm bước cho tới khi tìm được ô đầu tiên chứng
minh được, rồi dừng ngay. API dùng phương thức này cho chức năng gợi ý
(\texttt{GET /api/game/hint}), nên gợi ý không làm thay đổi ván chơi.
\texttt{classify\_all(state)} cũng chạy các bước đó nhưng không dừng sớm, và
trả về nhãn của \emph{mọi} ô chưa chứng minh; nó cho thấy toàn bộ những gì
$\mathit{KB}_t$ suy ra được tại một thời điểm. Cuối cùng,
\texttt{solve(engine)} lặp \texttt{next\_move} và nộp phán quyết cho tới khi
ván kết thúc. API dùng phương thức này cho chức năng tự giải
(\texttt{POST /api/game/auto-solve}).

\paragraph{Tác nhân không được tin tưởng mặc định.}
Đường ống có hai lớp kiểm tra độc lập. Tác nhân chỉ nộp phán quyết mà nó đã tự
chứng minh, nhưng Game Engine không dựa vào điều đó: \texttt{submit\_verdict}
chạy lại hai truy vấn trên cùng $\mathit{KB}_t$ và chỉ ghi nhận khi kết quả là
\texttt{ACCEPTED} (Bảng~\ref{tab:submit_outcome}). Cách bố trí này tốn thêm hai
lời gọi bộ giải cho mỗi phán quyết. Trên bản đồ Khu chợ chẳng hạn, tác nhân gọi
bộ giải $29$ lần và Game Engine gọi thêm $16$ lần, tổng cộng $45$ lần
(Bảng~\ref{tab:benchmark_raw}). Đổi lại, một lỗi trong tác nhân không thể đưa
ván chơi vào trạng thái sai: phán quyết không chứng minh được sẽ bị từ chối như
khi người chơi nộp.
\end{sloppypar}
